\chapter{Fourier Series}
\renewcommand{\currentchapterpath}{chapters/fs/}

\begin{flushright}
\textit{“八音克谐，无相夺伦。” ——《尚书·舜典》}
\footnote{From the ``Canon of Shun'' in the \textit{Book of Documents}: ``The eight classes of musical sounds achieve harmony, with none displacing the proper order of another.'' The line praises harmony formed from distinct components. Fourier analysis similarly represents a complicated signal as a superposition of elementary oscillations, each with its own frequency and amplitude; their ordered combination reconstructs the original behavior.}
\end{flushright}

Fourier series resolve periodic functions into elementary oscillations. They occur throughout signal processing, acoustics, optics, and the solution of differential equations on finite intervals. This chapter develops the orthogonal-function setting and the Fourier theorem for periodic functions~\cite{fourier,bracewell,lighthill}. The passage from a discrete trigonometric expansion to an integral transform is taken up in the next chapter.

\chapterinput{fourier_series}
\chapterinput{exercises}
